% ============================================================
%  TN1 — Mise en place du projet Veridian
%  INF1430 — Projet de fin d'études
%  TÉLUQ
%
%  Ce document décrit l'état initial, réel et fonctionnel
%  de l'infrastructure du projet au jour 1, avant tout
%  développement applicatif.
% ============================================================
\documentclass[12pt, a4paper]{article}

% ── Encodage et langue ──────────────────────────────────────
\usepackage[utf8]{inputenc}
\usepackage[T1]{fontenc}
\usepackage[french]{babel}

% ── Mise en page ────────────────────────────────────────────
\usepackage[top=2.5cm, bottom=2.5cm, left=3cm, right=2.5cm]{geometry}
\usepackage{setspace}
\onehalfspacing

% ── Typographie ─────────────────────────────────────────────
\usepackage{lmodern}
\usepackage{microtype}

% ── Liens et PDF ────────────────────────────────────────────
\usepackage[hidelinks, pdftitle={TN1 - Veridian}, pdfauthor={Francis Côté}]{hyperref}

% ── Tableaux et listes ──────────────────────────────────────
\usepackage{booktabs}
\usepackage{tabularx}
\usepackage{enumitem}
\usepackage{longtable}
\usepackage{array}

% ── En-têtes et pieds de page ───────────────────────────────
\usepackage{fancyhdr}
\pagestyle{fancy}
\fancyhf{}
\fancyhead[L]{\small INF1430 — Projet de fin d'études}
\fancyhead[R]{\small Francis Côté}
\fancyfoot[C]{\thepage}
\renewcommand{\headrulewidth}{0.4pt}

% ── Code source ─────────────────────────────────────────────
\usepackage{listings}
\usepackage{xcolor}

\definecolor{codebg}{RGB}{245,245,245}
\definecolor{codeframe}{RGB}{200,200,200}
\definecolor{codecomment}{RGB}{100,100,100}
\definecolor{codekeyword}{RGB}{0,100,180}
\definecolor{codestring}{RGB}{160,60,0}

% Style de listing uniforme — footnotesize pour éviter le débordement
% des lignes longues (commentaires inclus) dans la largeur de page.
\lstdefinestyle{veridian}{
  backgroundcolor=\color{codebg},
  frame=single,
  rulecolor=\color{codeframe},
  basicstyle=\ttfamily\footnotesize,
  keywordstyle=\color{codekeyword}\bfseries,
  stringstyle=\color{codestring},
  commentstyle=\color{codecomment}\itshape,
  breaklines=true,
  showstringspaces=false,
  numbers=left,
  numberstyle=\tiny\color{gray},
  numbersep=6pt,
  tabsize=2,
}
\lstset{style=veridian}

% ── Graphiques ──────────────────────────────────────────────
\usepackage{graphicx}
\usepackage{float}

% ── Références bibliographiques ─────────────────────────────
\usepackage[backend=biber, style=apa, sorting=nyt]{biblatex}
\addbibresource{references.bib}

% ── Largeurs de colonnes pour les longtables ────────────────
% Calculées pour une zone de texte de ~15.5cm (A4, marges 3+2.5cm).
% col1 = numéro bloc, col2 = titre, col3 = description/critères.
\newlength{\colbloc}   \setlength{\colbloc}{1.0cm}
\newlength{\coltitre}  \setlength{\coltitre}{4.8cm}
\newlength{\coldesc}   \setlength{\coldesc}{9.2cm}  % 15.5 - 1.0 - 4.8 - 0.5 (tabcol)

% ============================================================
%  DÉBUT DU DOCUMENT
% ============================================================
\begin{document}

% ────────────────────────────────────────────────────────────
%  PAGE TITRE
% ────────────────────────────────────────────────────────────
\begin{titlepage}
  \centering
  \vspace*{2cm}

  {\Large \textbf{Université TÉLUQ}}\\[0.5cm]
  {\large INF1430 — Projet de fin d'études}\\[3cm]

  \rule{\linewidth}{0.5pt}\\[0.4cm]
  {\LARGE \textbf{Travail noté 1}}\\[0.3cm]
  {\Large \textit{Mise en place du projet Veridian}}\\[0.3cm]
  \rule{\linewidth}{0.5pt}\\[3cm]

  \begin{tabular}{rl}
    \textbf{Étudiant :}    & Francis Côté \\[0.3cm]
    \textbf{Cours :}       & INF1430 — Projet de fin d'études \\[0.3cm]
    \textbf{Institution :} & Université TÉLUQ \\[0.3cm]
  \end{tabular}

  \vfill
  {\small Remis conformément aux exigences du cours INF1430}
\end{titlepage}

% ────────────────────────────────────────────────────────────
%  TABLE DES MATIÈRES
% ────────────────────────────────────────────────────────────
\newpage
\tableofcontents
\newpage

% ============================================================
\section{Dépôt Git}
% ============================================================

% Le dépôt est la première infrastructure mise en place.
% Il centralise le code, les migrations SQL, la documentation
% et l'historique de toutes les décisions techniques du projet.
Le dépôt public du projet est hébergé sur GitHub à l'adresse
\url{https://github.com/frankydoody/Veridian}. Il est créé
dès le premier jour du projet et constitue la source de vérité
unique pour l'ensemble du code source, des migrations de base
de données, de la documentation académique et de l'historique
des décisions techniques.

% ── 1.1 Structure des répertoires ───────────────────────────
\subsection{Structure des répertoires}

% La structure est définie au jour 1 et figée avant tout
% développement. Elle impose la séparation des responsabilités
% entre le backend, le frontend et la documentation.
La structure du dépôt est établie avant tout développement
applicatif. Elle impose une séparation stricte entre le serveur
Express, l'interface React et la documentation académique,
et restera stable pour toute la durée du projet.

\begin{lstlisting}[language=bash, caption={Arborescence planifiée du dépôt Veridian}]
Veridian/
├── start.ps1              # Demarrage automatise (Windows)
├── docker-compose.yml     # PostgreSQL + pgvector
├── Projet/
│   ├── Backend/           # API Node.js / Express 5
│   │   ├── app.js         # Config Express
│   │   ├── server.js      # Point d'entree HTTP
│   │   ├── package.json   # Dependances Node.js
│   │   ├── database/
│   │   │   └── migrations/  # SQL versiones
│   │   └── src/
│   │       ├── config/      # Connexion DB et IA
│   │       ├── controllers/ # Logique metier
│   │       ├── middlewares/ # Auth JWT, erreurs
│   │       ├── models/      # Requetes SQL
│   │       ├── routes/      # Cablage HTTP
│   │       ├── services/    # Whisper, Gemini, RAG
│   │       └── utils/       # Utilitaires communs
│   └── Frontend/          # React 19 / Vite
│       ├── index.html     # Point d'entree HTML
│       ├── package.json   # Dependances React
│       └── src/
│           ├── api/       # Client HTTP
│           ├── contexts/  # AuthContext
│           ├── components/# Composants reutilisables
│           ├── pages/     # Une page par route
│           └── utils/     # Fonctions partagees
└── Documentations/        # Rapports TN1 a TN4

\end{lstlisting}

% ── 1.2 Stratégie de branches ───────────────────────────────
\subsection{Stratégie de branches}

% La stratégie de branches est définie dès le départ et
% s'applique sans exception à chaque bloc fonctionnel.
% Les branches ne sont jamais supprimées : règle de
% traçabilité académique, pas une option.
La gestion du code source repose sur un modèle de branches
thématiques (\textit{feature branch workflow}). Chaque bloc
fonctionnel possède sa propre branche nommée \texttt{bloc-XX}.
Les règles suivantes s'appliquent sans exception :

\begin{enumerate}
  \item La branche est créée depuis \texttt{main} au démarrage du bloc.
  \item Le développement progresse jusqu'à un état fonctionnel et testable.
  \item Une \textit{pull request} est ouverte vers \texttt{main}.
  \item La fusion s'effectue exclusivement par \textit{merge commit}
        (\texttt{gh pr merge --merge}), sans \textit{squash} ni
        \textit{rebase}, afin de préserver l'intégralité de l'historique.
  \item La branche source \textbf{n'est jamais supprimée} après la fusion,
        conformément à la politique de traçabilité du projet.
  \item \texttt{main} demeure en état déployable à tout moment.
\end{enumerate}

% Tableau des branches — tabularx avec largeurs contrôlées
% pour éviter tout débordement sur la marge droite.
\begin{table}[H]
  \centering
  \caption{Branches planifiées pour le projet Veridian}
  \begin{tabularx}{\textwidth}{>{\ttfamily}l l X}
    \toprule
    \textbf{Branche}  & \textbf{Statut}  & \textbf{Contenu prévu} \\
    \midrule
    main      & Active (base) & Code de production intégré — toujours en état déployable \\
    bloc-01   & Planifiée     & Infrastructure et configuration initiale du projet \\
    bloc-02   & Planifiée     & Authentification JWT : inscription, connexion, middleware \\
    bloc-03   & Planifiée     & Gestion des projets CRUD \\
    bloc-04   & Planifiée     & Gestion des réunions CRUD \\
    bloc-05   & Planifiée     & Transcription audio avec Whisper (modèle local) \\
    bloc-06   & Planifiée     & Extraction des décisions via Gemini 1.5 Flash \\
    bloc-07   & Planifiée     & Mémoire vectorielle RAG avec pgvector \\
    bloc-08   & Planifiée     & Détection de contradictions et alertes \\
    bloc-09   & Planifiée     & Chat mémoire et export PDF \\
    bloc-10   & Planifiée     & Frontend Auth (React / Vite) \\
    bloc-11   & Planifiée     & Frontend Dashboard (Projets, Réunions, Tâches, Mémoire) \\
    bloc-12   & Planifiée     & Frontend connexion à l'API réelle \\
    bloc-13   & Planifiée     & Frontend Formulaires (Projet, Réunion, Tâche) \\
    bloc-14   & Planifiée     & Frontend Réunion Détaillée (audio, transcription, décisions) \\
    bloc-15   & Planifiée     & Frontend Décisions et Contradictions \\
    bloc-16   & Planifiée     & Frontend Calendrier de projet \\
    bloc-17   & Planifiée     & Tests d'intégration et qualité (Jest) \\
    bloc-18   & Planifiée     & Déploiement et documentation finale (TN4) \\
    \bottomrule
  \end{tabularx}
\end{table}

% ============================================================
\section{Board GitHub Projects}
% ============================================================

Un tableau de suivi est configuré dans l'onglet \textit{Projects}
du dépôt dès le démarrage du projet. Il centralise le suivi
de tous les blocs de développement et lie chaque carte à
l'issue GitHub et à la \textit{pull request} correspondantes.

% ── 2.1 Configuration de la vue ─────────────────────────────
\subsection{Configuration de la vue}

Le tableau utilise une vue \textit{Table} avec les champs suivants :

\begin{table}[H]
  \centering
  \caption{Champs de la vue Table — GitHub Projects}
  \begin{tabularx}{\textwidth}{lX}
    \toprule
    \textbf{Champ}                  & \textbf{Rôle} \\
    \midrule
    \textit{Title}                  & Titre du bloc et numéro d'issue associée \\
    \textit{Assignees}              & Responsable du bloc (frankydoody) \\
    \textit{Status}                 & État courant : \textit{Todo} ou \textit{Done} \\
    \textit{Linked pull requests}   & PR fusionnée liée automatiquement à l'issue \\
    \textit{Sub-issues progress}    & Progression des sous-tâches si applicables \\
    \bottomrule
  \end{tabularx}
\end{table}

% ── 2.2 État initial du tableau ─────────────────────────────
\subsection{État initial du tableau}

Au démarrage du projet, les dix-huit issues (blocs 01 à 18) sont
toutes en statut \textit{Todo}. Le statut passe automatiquement
à \textit{Done} lors de la fusion de la \textit{pull request}
associée dans \texttt{main}.
% ============================================================
\section{Backlog initial d'issues}
% ============================================================

% Toutes les issues sont créées au jour 1 via la CLI GitHub
% avant tout développement. Chaque issue décrit un bloc de
% travail autonome avec ses critères de complétion explicites.
Toutes les issues sont créées dès le premier jour via l'interface
en ligne de commande GitHub (\texttt{gh issue create}), avant
tout développement applicatif. Chaque issue représente un bloc
de travail autonome et livrable, portant un titre, une
description et des critères de complétion explicites.

% longtable avec largeurs p{} explicites : col1=numéro (1cm),
% col2=titre (4.8cm), col3=critères (reste ~9.2cm).
% Cela garantit que les trois colonnes tiennent dans la page.
\begin{longtable}{>{\centering}p{\colbloc} p{\coltitre} p{\coldesc}}
  \caption{Backlog initial — issues GitHub (état au jour 1)} \label{tab:backlog}\\
  \toprule
  \textbf{Bloc} & \textbf{Titre} & \textbf{Critères de complétion} \\
  \midrule
  \endfirsthead
  \multicolumn{3}{l}{\small\itshape Suite du tableau \ref{tab:backlog}…}\\[2pt]
  \toprule
  \textbf{Bloc} & \textbf{Titre} & \textbf{Critères de complétion} \\
  \midrule
  \endhead
  \midrule
  \multicolumn{3}{r}{\small\itshape Suite à la page suivante…}\\
  \endfoot
  \bottomrule
  \endlastfoot
  % ── Blocs infrastructure et backend ────────────────────────
  01 & Infrastructure et Configuration  & Dépôt créé, Docker opérationnel, migrations appliquées, \texttt{start.ps1} fonctionnel \\[3pt]
  02 & Authentification JWT             & Inscription, connexion, \texttt{authMiddleware}, token JWT retourné et validé \\[3pt]
  03 & Gestion des Projets CRUD         & Endpoints GET / POST / PUT / DELETE \texttt{/api/projects}, soft delete \\[3pt]
  04 & Gestion des Réunions CRUD        & Endpoints réunions, upload Multer, fichier audio persisté sur disque \\[3pt]
  05 & Transcription Audio (Whisper)    & Pipeline Whisper local (modèle \textit{small}), transcription persistée en base \\[3pt]
  06 & Extraction des décisions (Gemini)& Appel Gemini 1.5 Flash, décisions extraites et liées à la réunion \\[3pt]
  07 & Mémoire vectorielle (RAG)        & Embeddings générés, indexés dans pgvector, recherche sémantique opérationnelle \\[3pt]
  08 & Détection contradictions \& Alertes & Comparaison des décisions, alertes de contradiction affichées \\[3pt]
  09 & Chat mémoire \& Export PDF       & Interface de chat sur la mémoire RAG, export PDF des décisions \\[3pt]
  % ── Blocs frontend ─────────────────────────────────────────
  10 & Frontend Auth (React/Vite)       & Pages connexion et inscription, gestion du token JWT côté client \\[3pt]
  11 & Frontend Dashboard               & Pages Projets, Réunions, Tâches, Mémoire avec données statiques \\[3pt]
  12 & Frontend API Connect             & Toutes les pages connectées à l'API réelle, données dynamiques \\[3pt]
  13 & Frontend Formulaires             & Formulaires de création projet, réunion, tâche avec validation côté client \\[3pt]
  14 & Frontend Réunion Détaillée       & Lecteur audio intégré, transcription complète, décisions liées \\[3pt]
  15 & Frontend Décisions \& Contradictions & Vue agrégée des décisions, alertes en temps réel \\[3pt]
  16 & Frontend Calendrier de projet    & Vue calendrier mensuel, navigation entre les mois \\[3pt]
  % ── Blocs qualité et livraison ─────────────────────────────
  17 & Tests d'intégration \& Qualité   & Suite Jest sur les routes critiques : auth, transcription, mémoire RAG \\[3pt]
  18 & Déploiement \& Documentation     & Guide de déploiement, variables documentées, rapport TN4 \\
\end{longtable}

% ============================================================
\section{Plan de projet}
% ============================================================

% Ce plan décrit l'infrastructure initiale réelle et fonctionnelle
% mise en place au jour 1. Chaque élément est constatable
% directement dans le dépôt GitHub.
Le plan de projet est documenté dans \texttt{README.md} à la
racine du dépôt, complété par le présent document. Il décrit
l'infrastructure initiale réelle et fonctionnelle mise en
place dès le premier jour — chaque élément est constatable
directement dans le dépôt GitHub.

% ── 4.1 Découpage en blocs fonctionnels ─────────────────────
\subsection{Découpage en blocs fonctionnels}

% Les blocs sont séquentiels et dépendants : un bloc ne
% démarre qu'une fois le précédent fusionné dans main.
Le projet est structuré en dix-huit blocs fonctionnels séquentiels,
numérotés de 01 à 18. Chaque bloc produit un incrément livrable,
testé et intégré dans \texttt{main} avant que le suivant ne
commence. Ce séquençage garantit que l'application reste
déployable et démontrable à chaque étape du développement.

% Même logique de colonnes p{} que pour le backlog :
% col1=numéro, col2=nom court, col3=livrable attendu.
\begin{longtable}{>{\centering}p{\colbloc} p{\coltitre} p{\coldesc}}
  \caption{Plan de développement — blocs et livrables attendus} \label{tab:plan}\\
  \toprule
  \textbf{Bloc} & \textbf{Nom} & \textbf{Livrable attendu} \\
  \midrule
  \endfirsthead
  \multicolumn{3}{l}{\small\itshape Suite du tableau \ref{tab:plan}…}\\[2pt]
  \toprule
  \textbf{Bloc} & \textbf{Nom} & \textbf{Livrable attendu} \\
  \midrule
  \endhead
  \midrule
  \multicolumn{3}{r}{\small\itshape Suite à la page suivante…}\\
  \endfoot
  \bottomrule
  \endlastfoot
  % ── Blocs infrastructure et backend ────────────────────────
  01 & Infrastructure et Config   & Dépôt opérationnel, Docker configuré, migrations prêtes, \texttt{start.ps1} fonctionnel \\[3pt]
  02 & Authentification JWT        & API auth complète : inscription, connexion, middleware de protection des routes \\[3pt]
  03 & Projets CRUD                & Gestion complète des projets : création, lecture, mise à jour, suppression logique \\[3pt]
  04 & Réunions CRUD               & Création de réunions, upload et persistance du fichier audio via Multer \\[3pt]
  05 & Transcription Whisper       & Pipeline local Whisper (modèle \textit{small}) : transcription audio persistée en base \\[3pt]
  06 & Décisions Gemini            & Extraction automatique des décisions via Gemini 1.5 Flash, liées à chaque réunion \\[3pt]
  07 & Mémoire RAG                 & Embeddings vectoriels dans pgvector, recherche sémantique sur la mémoire de réunion \\[3pt]
  08 & Contradictions \& Alertes   & Détection automatique des contradictions entre décisions, alertes persistées \\[3pt]
  09 & Chat mémoire \& PDF         & Interface de chat sur la mémoire RAG, export PDF des décisions par projet \\[3pt]
  % ── Blocs frontend ─────────────────────────────────────────
  10 & Frontend Auth               & Pages connexion / inscription React, gestion du token JWT en contexte global \\[3pt]
  11 & Frontend Dashboard          & Toutes les pages avec données statiques (Projets, Réunions, Tâches, Mémoire) \\[3pt]
  12 & Frontend API Connect        & Pages connectées à l'API réelle : toutes les données dynamiques affichées \\[3pt]
  13 & Frontend Formulaires        & Formulaires de création et d'édition avec validation côté client \\[3pt]
  14 & Frontend Réunion Détaillée  & Lecteur audio intégré, transcription complète, décisions liées \\[3pt]
  15 & Frontend Décisions          & Vue agrégée des décisions par projet, alertes de contradiction en temps réel \\[3pt]
  16 & Frontend Calendrier         & Vue calendrier mensuel des réunions par projet, navigation entre les mois \\[3pt]
  % ── Blocs qualité et livraison ─────────────────────────────
  17 & Tests \& Qualité            & Suite Jest couvrant les routes critiques : auth, transcription, mémoire RAG \\[3pt]
  18 & Déploiement \& Docs         & Guide de déploiement, variables documentées, rapport final TN4 \\
\end{longtable}

% ── 4.2 Infrastructure initiale mise en place ───────────────
\subsection{Infrastructure initiale mise en place}

% Cette sous-section décrit ce qui est en place dès le jour 1,
% avant tout développement des blocs 01 à 18.
L'infrastructure suivante est mise en place dès le premier jour,
avant tout développement applicatif. Elle constitue le socle
technique sur lequel repose l'ensemble des blocs de développement.

\begin{itemize}
  % ── Dépôt et versionnement ──────────────────────────────────
  \item \textbf{Dépôt Git} : dépôt public sur GitHub, structure
        de répertoires établie, fichier \texttt{.gitignore} configuré,
        \texttt{README.md} documentant le projet et la procédure
        de démarrage.

  % ── Environnement conteneurisé ─────────────────────────────
  \item \textbf{Base de données} : fichier \texttt{docker-compose.yml}
        configuré avec l'image \texttt{pgvector/pgvector:pg16},
        incluant un \textit{healthcheck} et la persistance des données.
        Les fichiers de migration SQL (\texttt{001\_schema.sql},
        \texttt{002\_memory.sql}, \texttt{003\_tasks.sql}) sont
        versionnés dans le dépôt, prêts à être appliqués au
        premier démarrage.

  % ── Script de démarrage ─────────────────────────────────────
  \item \textbf{Script de démarrage} : \texttt{start.ps1} automatise
        la séquence complète sur Windows — vérification de Docker,
        démarrage du conteneur, application des migrations, installation
        des dépendances npm, ouverture des terminaux backend et frontend,
        puis lancement du navigateur. Le projet est opérationnel
        après un simple \texttt{git clone}.

  % ── Squelette backend ───────────────────────────────────────
  \item \textbf{Squelette backend} : structure de dossiers
        \texttt{src/} créée (config, controllers, middlewares,
        models, routes, services, utils), fichiers \texttt{app.js}
        et \texttt{server.js} initialisés, \texttt{package.json}
        avec toutes les dépendances définies.

  % ── Squelette frontend ──────────────────────────────────────
  \item \textbf{Squelette frontend} : projet Vite initialisé,
        React Router configuré, Tailwind CSS installé, client
        HTTP centralisé (\texttt{client.js}) et \texttt{AuthContext}
        en place.

  % ── Gestion de projet ───────────────────────────────────────
  \item \textbf{Gestion de projet} : 18 issues créées dans GitHub
        Projects, tableau Kanban configuré avec quatre colonnes,
        toutes les issues dans la colonne \textit{Backlog}.
\end{itemize}

% ── 4.3 Environnement d'exécution ───────────────────────────
\subsection{Environnement d'exécution}

% L'environnement est conteneurisé pour la base de données
% afin de garantir la reproductibilité sur tout poste
% disposant de Docker, indépendamment de l'OS.
L'environnement de développement est conteneurisé pour la
base de données, ce qui garantit la reproductibilité sur
n'importe quelle machine disposant de Docker.

\begin{lstlisting}[language=bash, caption={Fichier docker-compose.yml — configuration initiale}]
services:
  postgres:
    image: pgvector/pgvector:pg16   # PostgreSQL 16 + extension vectorielle
    container_name: veridian_db
    environment:
      POSTGRES_USER: veridian_user
      POSTGRES_PASSWORD: veridian_pass
      POSTGRES_DB: veridian
    ports:
      - "5432:5432"
    healthcheck:                    # Le backend attend que la DB soit prête
      test: ["CMD-SHELL", "pg_isready -U veridian_user -d veridian"]
      interval: 5s
      timeout: 5s
      retries: 10
\end{lstlisting}

L'image \texttt{pgvector/pgvector:pg16} est sélectionnée
spécifiquement pour son support natif de l'extension
\texttt{pgvector}, nécessaire au stockage des vecteurs
d'embeddings utilisés par le système de mémoire RAG.

% ── 4.4 Critères de complétion d'un bloc ────────────────────
\subsection{Critères de complétion d'un bloc}

% Ces quatre critères s'appliquent uniformément à chaque bloc.
% Un bloc n'est déclaré terminé que lorsque les quatre
% conditions sont simultanément satisfaites.
Un bloc est considéré \textbf{terminé} lorsque les quatre
critères suivants sont simultanément satisfaits :

\begin{enumerate}
  \item Le code est fusionné dans \texttt{main} via une
        \textit{pull request} approuvée (\texttt{merge commit}).
  \item L'application reste fonctionnelle et déployable
        après la fusion.
  \item L'issue GitHub correspondante est fermée et liée
        à la \textit{pull request}.
  \item La fonctionnalité est testable manuellement via
        l'interface ou directement via les endpoints de l'API.
\end{enumerate}

% ============================================================
%  RÉFÉRENCES
% ============================================================
\newpage
\printbibliography[title={Références}]

\end{document}
